% =============================================================================
% MOCK FINAL EXAM 3 (Simulado 3) --- Macroeconomics I, MPE Insper, 2026/T3
% Conceptual mock: every item asks to explain, draw, sign or say what an observer gets wrong.
% Exam only. The solution is a set of linked notes in Map/simulado-03/ (start at 00-index.md).
% Build from Simulados/:
%   pdflatex -jobname=simulado-03-exam simulado-03.tex     (run twice)
% Check fonts: pdffonts simulado-03-exam.pdf  (all Type 1).
% Every number in the solution notes: python simulado_codigo/check_simulado_03.py
% =============================================================================
\documentclass[11pt,a4paper]{article}
\usepackage[utf8]{inputenc}
\usepackage[T1]{fontenc}
\usepackage{lmodern}   % Type 1 vector fonts (replaces PK bitmaps)
\usepackage{cmap}      % ToUnicode CMap -> copyable text and accents
\usepackage[english]{babel}
\usepackage{amsmath,amssymb}
\usepackage[a4paper,margin=2.2cm]{geometry}
\usepackage{enumitem}
\usepackage{fancyhdr}
\usepackage{xcolor}
\usepackage{microtype}
\usepackage{booktabs}
\usepackage{needspace}

\newcommand{\disciplina}{Macroeconomics I --- MPE Insper}
\newcommand{\listatema}{Mock Final Exam 3 (conceptual)}

\pagestyle{fancy}\fancyhf{}
\lhead{\footnotesize\disciplina}
\rhead{\footnotesize \listatema}
\cfoot{\footnotesize\thepage}
\renewcommand{\headrulewidth}{0.3pt}
\setlength{\headheight}{14pt}
\setlist{topsep=2pt,itemsep=2pt}

\newcounter{questao}
\newcommand{\bloco}[3]{\Needspace{10\baselineskip}\bigskip\noindent
  {\large\bfseries\color{blue!55!black}Block #1 --- #2}\hfill\textit{(#3)}\par
  \noindent\rule{\linewidth}{0.3pt}\smallskip}
\newcommand{\questao}[1]{\Needspace{16\baselineskip}\refstepcounter{questao}\medskip\noindent
  \textbf{Question \thequestao.}~#1\par}
% objects the answer must contain
\newcommand{\objetos}[1]{\par\smallskip\noindent{\small\color{black!65}\textit{Answer every object:} #1}\par}
\newcommand{\espaco}[1]{\par\textit{Answer:}\vspace{#1}\par}
\newcommand{\vf}[2]{\Needspace{9\baselineskip}\par\medskip\noindent\textbf{#1)} (5 points) #2
  \quad\textbf{T} $\square$\ \ \textbf{F} $\square$\par\noindent\textit{Justification:}\vspace{4.6cm}\par}

\begin{document}
\begin{center}
  {\Large\bfseries \disciplina}\\[2pt]
  {\large \listatema\ --- Simulado 3 --- 2026, 3rd term}
\end{center}
\vspace{2pt}

{\footnotesize\noindent\textbf{Instructions.} Four blocks of 25 points (100 in total). Duration: 3 hours,
closed book, non-programmable calculator allowed. Each block has one True/False question (two items,
5 points each) and one discursive question (items a, b, c, 5 points each).

\smallskip\noindent\emph{True/False with justification.} Mark T or F and justify. The verdict is worth
\textbf{1 point} and the justification \textbf{4 points}. A verdict with no justification earns at most 1
point; a justification whose words contradict its own verdict or its own algebra earns 0. A good
justification names the mechanism in one causal chain and, when the statement is false, says what is true
instead.

\smallskip\noindent\emph{Discursive items} ask you to \emph{explain}, \emph{draw and interpret},
\emph{sign and say which effect dominates}, or say \emph{what an observer with ordinary data would wrongly
conclude}. Each item ends with the list of objects your answer must contain: tick them off before moving
on. Computation is asked for only where a short derivation is the explanation. Sign every derivative before
interpreting it; box final answers; label axes, curves and the direction of every shift.
Notation: Kurlat (2020) in Blocks I--III, Benigno (2015) in Block IV. All numbers are illustrative.

\smallskip\noindent\emph{Solution.} There is no answer key in this PDF. The worked solution, with every step,
the figures and the grading rubric, is the set of linked notes in \texttt{Map/simulado-03/}, starting at
\texttt{00-index.md}. Sit the exam timed first, then read.\par}

% =============================================================================
\bloco{I}{Measurement and national accounts}{Lecture 1 \,\textbullet\, 25 points}

\questao{(10 points) True or False? Justify.}

\vf{a}{A small economy produces only streaming subscriptions and cinema tickets (prices in reais,
quantities in millions):
\begin{center}\small
\begin{tabular}{lcccc}\toprule
 & $p_{\text{stream}}$ & $q_{\text{stream}}$ & $p_{\text{cinema}}$ & $q_{\text{cinema}}$\\\midrule
Year 1 & 40 & 10 & 20 & 30\\
Year 2 & 20 & 30 & 25 & 25\\\bottomrule
\end{tabular}
\end{center}
\emph{Statement:} ``A statistician who values both years' quantities at year-2 prices will report faster
real growth than one who uses year-1 prices, because year-2 prices are more up to date.''}

\vf{b}{\emph{Statement:} ``In a year in which the only price that changes is the price of imported
televisions, which rises by 30\%, the GDP deflator rises by more than the consumer price index (IPCA).''}

\questao{(15 points) Brazil and Paraguay: which number answers which question?
Over the last decade (illustrative figures), Paraguay's real GDP grew 3.5\% a year and Brazil's 2.5\%
a year; Paraguay's population grew 2.0\% a year and Brazil's 0.5\%. Today GDP per capita converted at market
exchange rates is USD 9{,}000 in Brazil and USD 6{,}000 in Paraguay. The price level relative to the United
States (the ratio of the PPP exchange rate to the market rate) is 0.55 in Brazil and 0.40 in Paraguay.
Employment is 45\% of the population in Brazil and 50\% in Paraguay.}

\medskip\noindent\textbf{a)} (5 points) A newspaper headline reads: ``Paraguay grows 40\% faster than Brazil,
so Paraguayans' living standards are catching up with Brazilians'.'' Compute the growth rate of GDP per
capita in each country (exactly and by the log approximation). Does the headline's conclusion follow?
Explain why aggregate growth and per capita growth answer different questions.
\objetos{$g_y$ for Paraguay; $g_y$ for Brazil; verdict on the headline; the question each growth rate answers.}
\espaco{6.5cm}

\Needspace{12\baselineskip}\noindent\textbf{b)} (5 points) Compute GDP per capita at PPP in both countries, and the Brazil/Paraguay
ratio at market rates and at PPP. Explain, naming the mechanism, why the ratio shrinks, and say what an
observer who uses market exchange rates would wrongly conclude about the gap in living standards.
\objetos{two PPP values; two ratios; the mechanism (which goods, which prices, why); the wrong conclusion
and its direction.}
\espaco{6.5cm}

\Needspace{12\baselineskip}\noindent\textbf{c)} (5 points) A consultant ranks the two countries' \emph{labour productivity} by GDP per
capita at PPP. Compute GDP per worker at PPP. Say which of three measures (aggregate GDP at market rates;
GDP per capita at PPP; GDP per worker at PPP) answers each of: (i) the size of the market an exporter
selling in dollars faces; (ii) average material living standards; (iii) labour productivity. Explain why per
capita and per worker give different ratios here.
\objetos{two per-worker values and their ratio; one measure for each of (i), (ii), (iii); the reason the two
ratios differ.}
\espaco{6.5cm}

% =============================================================================
\bloco{II}{Long-run growth and the Solow model}{Lectures 2--3 \,\textbullet\, 25 points}

\questao{(10 points) True or False? Justify. Solow model, variables per worker, no technological progress.}

\vf{a}{\emph{Statement:} ``The Solow model predicts that poorer countries grow faster than richer
ones. So a poor country that has grown more slowly than a rich one for decades is evidence against the
model.''}

\vf{b}{Take $\alpha=1/3$, $s=0.2$, $\delta=0.05$. \emph{Statement:} ``A permanent fall in the population
growth rate from 2\% to 1\% a year lowers, in the long run, both the growth rate of aggregate GDP and the
growth rate of GDP per capita.''}

\questao{(15 points) Tax paperwork read as productivity.
Brazilian firms spend an unusually large number of hours complying with tax rules. Model it as follows.
Output is $Y=AK^{\alpha}L_P^{1-\alpha}$, where $L_P$ is production labour. For each production worker, a
firm must also employ $\varphi$ compliance workers, who produce nothing and are paid the same wage $w$.
Official statistics record total employment $L=(1+\varphi)L_P$. Factor markets are competitive and the
price of output is 1. Take $\alpha=1/3$ and $\varphi=0.25$.}

\medskip\noindent\textbf{a)} (5 points) Write $Y$ as a function of $A$, $K$, $L$ and $\varphi$. An observer
has accurate data on $Y$, $K$, $L$ and on factor incomes, but cannot tell compliance workers from
production workers. Which labour share and which productivity $\hat A$ does he measure? Is the labour share
distorted? Where does the distortion show up?
\objetos{$Y(A,K,L,\varphi)$; measured labour share, with the firm's first-order condition; $\hat A/A$
(number); share of the workforce in compliance; where the wedge sits.}
\espaco{7cm}

\Needspace{12\baselineskip}\noindent\textbf{b)} (5 points) A tax simplification lowers $\varphi$ from 0.25 to 0.10 over five years.
True $A$ does not change. Write the growth-accounting decomposition of $g_Y$, solve for the Solow residual
and compute its annual value. What does the statistical agency report, and why is this consistent with the
reading of the residual as a measure that also captures policies that change how well resources are
allocated? What would it report if compliance rules got \emph{heavier}?
\objetos{decomposition; residual in terms of $\varphi$; its annual value; what is reported and why; the
opposite case.}
\espaco{7cm}

\Needspace{12\baselineskip}\noindent\textbf{c)} (5 points) Now let the economy follow the Solow model with $s$, $n$ and $\delta$
constant, per worker of \emph{measured} employment. Draw the effect of the reform in two panels: the
mechanism (Solow diagram) and the time path of $y$. By how much does steady-state $y^*$ rise? Sign the growth
rate of $y$ on impact, during the transition and in the long run, and explain why the total rise in $y$
exceeds the jump in measured TFP.
\objetos{two-panel diagram; $y^{*\prime}/y^*$; growth of $y$ in three horizons; the source of the extra gain.}
\espaco{7cm}

% =============================================================================
\bloco{III}{Microfoundations: consumption, labour and general equilibrium}{Lectures 4--6 \,\textbullet\, 25 points}

\questao{(10 points) True or False? Justify.}

\vf{a}{A household has one unit of time, utility $\ln c+b\ln l$ over consumption and leisure, earns the
after-tax wage $(1-\tau)w$ per hour worked and receives a lump-sum transfer $T\geq0$.
\emph{Statement:} ``If the government returns nothing to households ($T=0$), a rise in the labour-income tax
$\tau$ reduces hours worked, because it makes leisure cheaper.''}

\vf{b}{Matches are $m=\mu\,v^{\alpha}u^{1-\alpha}$ (in rates) and a fraction $s$ of employed workers loses
their job each period. \emph{Statement:} ``A fall in the unemployment rate together with a rise in the
vacancy rate, with matching efficiency $\mu$ and the separation rate $s$ unchanged, is an outward shift of
the Beveridge curve: the market now has more vacancies per unemployed worker.''}

\questao{(15 points) Interest rates, a tax rebate and a tax on interest income.
Households live two periods, have utility $u(c_1)+\beta u(c_2)$ with CRRA $u$ (coefficient $\sigma$; log
when $\sigma=1$), and borrow and lend at $r$ unless a limit is stated.}

\medskip\noindent\textbf{a)} (5 points) A saver has $y_1=100$, $y_2=0$, $\beta=0.96$, $\sigma=2$. The Copom
raises the real rate from 4\% to 6\%. Sign $\partial c_1/\partial r$, say which effect dominates and why,
and compute $c_1$ at both rates. Then suppose \emph{all} households are identical and capital and labour are
fixed, so output is given in each period (frozen factors). Explain why aggregate saving must be zero, what
then determines $r$, and compute $r$ when $y_2=y_1$ and when $y_2/y_1=1.03$.
\objetos{sign; dominant effect in words that agree with the sign; $c_1$ at 4\% and 6\%; why aggregate saving
is zero; what sets $r$; the two values of $r$.}
\espaco{7cm}

\Needspace{12\baselineskip}\noindent\textbf{b)} (5 points) A young household has $y_1=40$, $y_2=84$, taxes $\tau_1=10$,
$\tau_2=10.5$, log utility, $\beta=1/1.05$ and $r=5\%$. The government cuts $\tau_1$ to 0 and raises $\tau_2$
to 21. (i) With free borrowing, compute $c_1$, $c_2$ and $a$ before and after. (ii) Now impose $a\ge -15$.
Compute $c_1$, $c_2$ and $a$ before and after. Does the timing of taxes matter in each case? Give two
kinds of household for which (ii) is the relevant case.
\objetos{$c_1,c_2,a$ in all four cases; verdict on timing in (i) and in (ii), each with its reason; two
examples.}
\espaco{7cm}

\Needspace{12\baselineskip}\noindent\textbf{c)} (5 points) A saver has $y_1=100$, $y_2=0$, log utility, $\beta=0.96$, $r=5\%$. The
government taxes the return to saving: $c_2=(1+r-\tau_a)a$ with $\tau_a=0.02$. Compare with a lump-sum tax
in period 2 that raises the same revenue. Compute $c_2/c_1$ under each tax and the revenue. Which margin
does each tax distort? Show that the household could afford the savings-tax plan under the lump-sum tax,
and conclude which tax is better for the household.
\objetos{revenue; both ratios $c_2/c_1$; the margin each tax distorts; the affordability argument; the
welfare ranking.}
\espaco{7cm}

% =============================================================================
\bloco{IV}{Money, inflation and the New-Keynesian AS--AD model}{Lectures 7--9 \,\textbullet\, 25 points}

\questao{(10 points) True or False? Justify.}

\vf{a}{\emph{Statement:} ``Last year M2 grew 10\%, inflation was 6\% and real GDP grew 3\%. Since money grew
faster than nominal GDP, these data refute the quantity equation $MV=PY$.''}

\vf{b}{Money demand follows Baumol--Tobin, the real interest rate is 2\% and real output does not grow.
\emph{Statement:} ``Because money is neutral, a permanent rise in money growth from 4\% to 10\% a year has
no effect on any real variable in the new steady state.''}

\questao{(15 points) The Copom, a drought and a payroll tax.
Use Benigno (2015): AS is $p-p^e=\kappa(y-y_n)$ and AD is $y=\bar y_n-\sigma\left[i-(\bar p-p)-\rho\right]$,
with $\bar p=p^e$. Take $\sigma=0.5$, $\kappa=1.133$, $\eta=0.2$ and $\theta=8$. Output and prices are in
percent deviations from the initial natural level and from $p^e$.}

\medskip\noindent\textbf{a)} (5 points) Money-market equilibrium is $M^S=p\,m^D(Y,i)$. Explain what this
condition determines (i) when the central bank sets $M$ and (ii) when it sets $i$. For an exogenous rise in
$M$, say what adjusts with prices fixed and with prices flexible. Draw the two regimes. Why does Benigno's
AS--AD model have no LM curve?
\objetos{what is endogenous in (i) and in (ii); fixed-price and flexible-price answers; the two-panel diagram;
the reason for no LM curve.}
\espaco{7cm}

\Needspace{12\baselineskip}\noindent\textbf{b)} (5 points) A drought temporarily lowers agricultural productivity so that $y_n$ falls
by 2\% ($\bar y_n$ unchanged). The Copom keeps $i$ unchanged. Which curve shifts, in which direction, and
which curve does the economy move along? Compute $dy$, $dp$, $y-y_n$ and $y-y_e$. What change in $i$ closes
both gaps at once, and why can it?
\objetos{the curve that shifts and its direction; the movement along the other; $dy$, $dp$; both gaps with
signs; the required $di$; why one instrument suffices.}
\espaco{7cm}

\Needspace{12\baselineskip}\noindent\textbf{c)} (5 points) Instead, a temporary payroll-tax rise raises the mark-up by $d\mu=4.4\%$,
so that $y_n$ falls by 2\% while $y_e$ does not move. The Copom minimises
$L=(y-y_e)^2+\frac{\theta}{\kappa}(p-p^e)^2$. Explain why the AS curve, and not the AD curve, is the
constraint of this problem. Derive the optimal point and the change in $i$ that delivers it, and explain why
there is a trade-off here but not in (b).
\objetos{why AS is the constraint; the targeting rule; $y-y_e$ and $p-p^e$ at the optimum; $di$; the
trade-off argument.}
\espaco{7cm}

\vfill
\begin{center}\footnotesize\textit{End of exam. Solution: \texttt{Map/simulado-03/00-index.md}.}\end{center}
\end{document}
